\documentclass[a4paper,12pt]{article}

%%% Работа с русским языком
%\usepackage{DejaVuSans}
\usepackage{cmbright}
%\usepackage[OT1]{fontenc}
\renewcommand*\familydefault{\sfdefault} %% Only if the base font of the document is to be sans serif
\usepackage{cmap}					% поиск в PDF
\usepackage{mathtext} 				% русские буквы в формулах
\usepackage[T2A]{fontenc}			% кодировка
\usepackage[utf8]{inputenc}			% кодировка исходного текста
\usepackage[english]{babel}	% локализация и переносы
\usepackage{indentfirst}            % красная строка в первом абзаце
\usepackage{mathtext} 				% русские буквы в формулах
\frenchspacing                      % равные пробелы между словами и предложениями



%Добавлено:
\usepackage{refcount} 				% перекрёстные ссылки внутри документа
\usepackage{fixfoot,xspace}			% двойная ссылка
\usepackage{tipa}					% греческие в тексте
\usepackage{dcolumn}				% выравнивание по точке
\usepackage{mhchem} 				% Оформление химических формул
%\usepackage{acro}					% ведение списка сокращений
%\AddAcroTranslations{Russian}{list-name = Список сокращений} 	% Что будет написано в заголовке
\usepackage{caption}				% несколько рисунков
\usepackage{subcaption}				% несколько рисунков
\usepackage{gensymb}				% градусы цельсия
\usepackage{multicol}				% Текст в две колонки
\usepackage{siunitx}				% работа с цифрами
\usepackage{wrapfig}				% обтекание рисунков текстом
\usepackage[inline]{enumitem}		%Горизонтальный список, добавлется звёздочкой

%Добавленные команды
\newcommand{\grad}{~\celsius}

% \makeatletter									%Подпись к рисункам Рисунок, вместо Рис
% \renewcommand{\fnum@figure}{Рисунок \thefigure}
% \makeatother
% %\counterwithin{figure}{section}

\renewcommand{\labelenumii}{\arabic{enumi}.\arabic{enumii}}	% отвечает за уровень списка
\renewcommand{\labelenumiii}{\arabic{enumi}.\arabic{enumii}.\arabic{enumiii}}
\renewcommand{\labelenumiv}{\arabic{enumi}.\arabic{enumii}.\arabic{enumiii}.\arabic{enumiv}}

% \addto\captionsrussian{% Replace "english" with the language you use
%   \renewcommand{\contentsname}{Table of Contents}}%%Подпись к оглавлению



%%% Шаблонная информация для титульного листа
\newcommand{\CourseName}{}
\newcommand{\FullCourseNameFirstPart}{\so{}}
\newcommand{\SemesterNumber}{}
\newcommand{\LecturerInitials}{}
\newcommand{\CourseDate}{2023}

%%% Оформление страницы
\usepackage{extsizes}     % Возможность сделать 14-й шрифт
\usepackage{geometry}     % Простой способ задавать поля
\usepackage{setspace}     % Интерлиньяж
\usepackage{enumitem}     % Настройка окружений itemize и enumerate
\setlist{leftmargin=25pt} % Отступы в itemize и enumerate

\geometry{top=15mm}    % Поля сверху страницы
\geometry{bottom=30mm} % Поля снизу страницы
\geometry{left=15mm}   % Поля слева страницы
\geometry{right=15mm}  % Поля справа страницы

\setlength\parindent{15pt}        % Устанавливает длину красной строки
\linespread{1.1}                  % Коэффициент межстрочного интервала
%\setlength{\parskip}{0.2em}      % Вертикальный интервал между абзацами
%\setcounter{secnumdepth}{0}      % Отключение нумерации разделов
%\setcounter{section}{-1}         % Нумерация секций с нуля
\usepackage{multicol}			  % Для текста в нескольких колонках
\usepackage{soulutf8}             % Модификаторы начертания

%%% Добавлено:
\usepackage{parskip} % сам настроит интервалы нужным образом
\setlength{\parindent}{0,8cm} % по умолчанию — 0

%%% Колонтитулы
\usepackage{titleps}
\newpagestyle{main}{
	\setheadrule{0pt}
%	\sethead{\subsectiontitle}{}{\hyperlink{intro}{\;К содержанию}}
	\setfootrule{0pt}                       
	\setfoot{}{}{\thepage} 
}
\pagestyle{main} 

%%% Дополнительная работа с математикой
\usepackage{amsmath,amsfonts,amssymb,amsthm,mathtools} % пакеты AMS
\usepackage{icomma}                                    % "Умная" запятая

%\newcommand{\dse}{\displaystyle}

%%% Перенос знаков в формулах (по Львовскому)
%\newcommand*{\hm}[1]{#1\nobreak\discretionary{}{\hbox{$\mathsurround=0pt #1$}}{}}

%%% Работа с картинками
\usepackage{graphicx}    % Для вставки рисунков
\setlength\fboxsep{3pt}  % Отступ рамки \fbox{} от рисунка
\setlength\fboxrule{1pt} % Толщина линий рамки \fbox{}
\usepackage{wrapfig}     % Обтекание рисунков текстом

%%% Работа с таблицами
\usepackage{array,tabularx,tabulary,booktabs} % Дополнительная работа с таблицами
\usepackage{longtable}                        % Длинные таблицы
\usepackage{multirow}                         % Слияние строк в таблице
\usepackage{makecell}						% Перенос в ячейках строк
%\usepackage{tablefootnote}					% Сноски в таблицах (работает плохо)
\usepackage{booktabs,caption}				% Сноски под таблицей
\usepackage[flushleft]{threeparttable}		% Сноски под таблицей


%%% Теоремы
\theoremstyle{plain}
\newtheorem{theorem}{Теорема}[section]
\newtheorem{lemma}{Лемма}[section]
\newtheorem{proposition}{Утверждение}[section]
\newtheorem*{exercise}{Упражнение}

\theoremstyle{definition}
\newtheorem{definition}{Определение}[section]
\newtheorem*{adefinition}{Определение(не материал лектора)}
\newtheorem*{corollary}{Следствие}
\newtheorem*{addition}{Дополнение}
\newtheorem*{note}{Замечание}
\newtheorem*{anote}{Замечание автора}
\newtheorem*{reminder}{Напоминание}
\newtheorem{example}{Пример}
\newtheorem*{examples}{Примеры}

\theoremstyle{remark}
\newtheorem*{solution}{Решение}


 

% %%% Нумерация уравнений
% \makeatletter
% \def\eqref{\@ifstar\@eqref\@@eqref}
% \def\@eqref#1{\textup{\tagform@{\ref*{#1}}}}
% \def\@@eqref#1{\textup{\tagform@{\ref{#1}}}}
% \makeatother                      % \eqref* без гиперссылки
% \numberwithin{equation}  % Нумерация вида (номер_секции).(номер_уравнения)
% \mathtoolsset{showonlyrefs=false} % Номера только у формул с \eqref{} в тексте.

%%% Гиперссылки
\usepackage{hyperref}
\usepackage[usenames,dvipsnames,svgnames,table,rgb]{xcolor}
\hypersetup{
	unicode=true,            % русские буквы в раздела PDF
	colorlinks=true,       	 % Цветные ссылки вместо ссылок в рамках
	linkcolor=black!15!blue, % Внутренние ссылки
	citecolor=green,         % Ссылки на библиографию
	filecolor=magenta,       % Ссылки на файлы
	urlcolor=NavyBlue,       % Ссылки на URL
}

%%% Графика
\usepackage{tikz}        % Графический пакет tikz
\usepackage{tikz-cd}     % Коммутативные диаграммы
\usepackage{tkz-euclide} % Геометрия
\usepackage{stackengine} % Многострочные тексты в картинках
\usetikzlibrary{angles, babel, quotes, automata,arrows,calc,positioning}

%%% Xlam
%\usepackage{longtable}
%\newcolumntype{d}[1]{D{.}{.}{#1}}
%\usepackage{paratype}
%\renewcommand{\thefootnote}{\alph{footnote}} % Делает сноски буквами


